\section{Lens specifications}
\label{app:lenses}

This appendix gives the seven methodology lenses (\cref{sec:foundations,sec:n3}) at
implementation grain: what each one fires on, what it treats as missing, what kind of knowledge
it predicts, how a human would be asked, and where it is known to fail. Source: \repo{gapmap/src/gapmap/lenses.py},
\repo{gapmap/src/gapmap/semantic.py} (closure judging) and \repo{docs/plans/n3-poc-lens-spec.md}
\S2 and \S2.8 (post-review text; \S8/\S9's numbers are the pre-review prototype's and are not used
here). A lens is an analytical construct applied uniformly to a ledger, never a named research
method by itself (\cref{sec:foundations}).

\paragraph{Shared machinery, once.} Every lens works over an \evobs{} \code{Ledger}: claims labeled
{\small\nolinkurl{literature_supported}} or {\small\nolinkurl{organisational_artefact_supported}} form pool
\textbf{A} (attestable); {\small\nolinkurl{synthetic_extrapolation}} claims form pool \textbf{S} (never attestation,
usable only to flag an unverified statement). A lens \emph{fires} on an eligible \textbf{A} claim
(the seed), producing a \code{Candidate}. Whether the candidate's missing element is genuinely
absent (\emph{closure}, \cref{sec:n3}) is decided downstream by a local, replayed semantic judge
(\code{qwen2.5:7b-instruct}), never by the lens itself; each lens also computes a comparison-only
\emph{lexical state} (regex-based) that is never decisive but feeds the \S7.1 checks
(\cref{app:claims}, rows E36, E38). A seed sourced from a promotional domain or matching the
\code{PROMO} lexicon never fires, for any lens (\code{promo_blocked}, \repo{gapmap/src/gapmap/lenses.py}).

\paragraph{Assertion, not always the span.} Every firing rule, every anchor, and the judge's own
probe text run on a claim's \code{assertion} -- the extractor's paraphrase of its source -- not
necessarily a verbatim sentence. A human-facing question quotes the most relevant verbatim span
sentence when one exists; \code{lenses.py}'s \code{_best_sentence} falls back to the assertion
itself when a claim has no span sentences, when the best candidate sentence opens with a bare
referent, or when it is shorter than five words. This is a construct-validity threat, not an edge
case: DIAG's anchor (below) is always built from the seed's assertion, never a span sentence, and
one trace in \cref{app:records} (Trace~3) quotes the model's paraphrase as if it were the
document's own words.

\subsection{DISC --- undefined discrimination}

\begin{description}
  \item[Construct and source] A cue is a discrimination: contrasting cases plus the feature that
    separates them. CTA interviews probe for cues, discriminations and
    typicality; \emph{which} cue triggered a judgment is behavior a person cannot fully narrate.
  \item[Firing rule] Fires on a term from the \code{JUDGE} lexicon (e.g.\ \enquote{appropriate},
    \enquote{high-risk}, \enquote{normal}) inside a verified claim's assertion. The anchor is that
    term plus the stem of the next content word (two fixed anchors, \enquote{high risk} and
    \enquote{large scale}, are exempt from this pairing). \textbf{Termhood:} an anchor needs at
    least 2 distinct \textbf{A} claims using it, or it is dropped (removes one-off bigrams such as
    \enquote{appropriate tank}).
  \item[Missing element] A boundary for the anchor: a positive/negative contrast, or a quantified
    threshold. The judge checks every claim whose text contains the anchor; a one-sided list of
    examples (\enquote{such as children}) counts only as \emph{partial}, never closed, because it
    gives no separating feature.
  \item[Predicted knowledge type] \code{cue} (\code{expectancy} if the term is \enquote{normal},
    \enquote{expected} or \enquote{baseline}). Tacitness: \code{perceptual} if at least half the
    anchor claims are typed failure-mode, cue, check, procedure-step or interpretation, otherwise
    \code{collective}.
  \item[Elicitation channel] Contrasting-case classification (perceptual knowledge), or
    expert-rated written vignettes (collective/relational knowledge) --- not an interview alone.
  \item[Question template (verbatim)] \enquote{Two sources write: \textquotedbl\{s1\}\textquotedbl{}
    and \textquotedbl\{s2\}\textquotedbl. Think of a recent case where you had to judge whether
    something was \{anchor\}. Describe one case that clearly was, one that clearly was not, and one
    that was hard to call. What differed between them, and what did you check first?}
  \item[Known failure modes] (i) the boundary is stated in paraphrase, which the anchor-string
    regex cannot see --- the main limitation of DISC; (ii) the term is
    \emph{institutionally indeterminate} (at least half the anchor claims are typed \code{norm} or
    \code{concept}), so practitioners may simply disagree and there is nothing to recover; (iii)
    the boundary sits in a source the reconstruction never fetched. Four of eight failures examined
    by hand were paraphrase misses (\repo{docs/plans/n3-poc-lens-spec.md} \S10).
\end{description}

\subsection{DIAG --- rival causes without a discriminating sign}

\begin{description}
  \item[Construct and source] PARI's Interpretation step updates a fault hypothesis among rival
    causes of one observation; Evidence-Centered Design's \enquote{additional KSAs} formalize the
    same rival-explanation structure.
  \item[Firing rule] Fires on an \textbf{A} claim typed \code{failure_mode}, \code{cue},
    \code{interpretation}, \code{misconception} or \code{expert_novice_contrast} whose assertion
    matches the \code{CAUSE} lexicon; the assertion is split at the marker into a cause and an
    effect side. Two cause claims are \emph{rivals} if their effect stems share $\geq 2$ stems and
    their non-common cause stems share none, and rivals are direct neighbors only (never a
    transitive chain). A seed plus its direct rivals is a group; a group needs $\geq 2$ rivals.
  \item[Missing element] A discriminating sign per rival cause. The local judge probes each rival
    in turn (its own cause/effect stems, matched against retrieved sentences via \code{DISCR});
    the group stays open if at least two rivals get no sign.
  \item[Predicted knowledge type] \code{interpretation} plus \code{cue}. Tacitness:
    \code{relational}, with \code{perceptual} added if any seed is typed \code{cue}.
  \item[Elicitation channel] CDM probes on a recalled case, then contrasting cases.
  \item[Question template (verbatim, built after judging)] \enquote{Sources list several causes of
    \{anchor\}: \textquotedbl\{s1\}\textquotedbl, \textquotedbl\{s2\}\textquotedbl. Think of the
    last time you diagnosed \{anchor\}. Which cause did you suspect first, and what did you see,
    hear or measure that let you rule the others out?} (\repo{gapmap/src/gapmap/semantic.py}:268--273;
    the lead sentence drops to one quote, or none, if fewer are available.) \textbf{The template
    substitutes \code{cand.anchor}, not a clean \enquote{effect} noun phrase:}
    \code{_diag_anchor()} (\repo{gapmap/src/gapmap/lenses.py}:364--375) returns the full sentence
    containing the cause marker, trimmed to 80 characters, or the whole seed assertion if the
    effect side is too short -- the same fragment used as the DIAG anchor everywhere else in this
    record. This is a known template defect (the rendered question can read as a mid-sentence
    fragment; \cref{app:defects}, row TR1), not a design choice.
  \item[Known failure modes] (i) the sign exists but is phrased without the cause's own stems ---
    lexical linking, not the sign's absence, is the binding constraint; (ii) the causes may
    co-occur, so practitioners test rather than discriminate between mutually exclusive options;
    (iii) the \enquote{sign} DISCR finds may be instrument output (a displayed reading), not tacit
    judgment. A first build of this lens read \enquote{rivals} as a transitive closure over shared
    stems instead of direct neighbors, chaining unrelated causes into one 40-claim group and
    inflating that group's rank; caught by reading the rendered map, not by the (all-passing) unit
    tests (\repo{docs/lessons.md}).
\end{description}

\subsection{SEL --- selection criterion named, mapping absent}

\begin{description}
  \item[Construct and source] Conditions for choosing between methods: GOMS selection rules, HTA
    plans, and the condition part of a knowledge component; \enquote{when (not) to apply} a rule.
  \item[Firing rule] Fires on an \textbf{A} assertion that matches the \code{SEL} lexicon
    (\enquote{based on}, \enquote{depending on}, \enquote{choose}, \ldots) \emph{and} is
    prescriptive (a \code{MODAL} marker, or typed \code{decision}).
  \item[Missing element] A mapping from situation to option: a conditional (\code{COND}) linking a
    circumstance to the named choice.
  \item[Predicted knowledge type] \code{decision}; tacitness \code{relational}.
  \item[Elicitation channel] CDM probes for the options considered and the basis of the choice.
  \item[Question template (verbatim)] \enquote{\textquotedbl\{s1\}\textquotedbl. Think of the last
    time you made this choice. Which options did you consider, which did you take, and what about
    that situation decided it? When did you last choose differently?}
  \item[Known failure modes] (i) the mapping is stated as an example or a comparative that
    \code{COND} does not match; (ii) the choice named is not a real decision point in field
    practice at all --- a prescriptive sentence need not reflect a live decision.
\end{description}

\subsection{HEDGE --- hedged rule, exception condition absent}

\begin{description}
  \item[Construct and source] \enquote{When \emph{not} to apply} a rule: case-based-reasoning
    relevance conditions and the condition part of a Knowledge-Learning-Instruction unit. A
    quasi-universal hedge (\enquote{generally}, \enquote{as a rule}) admits exceptions without
    stating them.
  \item[Firing rule] Fires on a prescriptive \textbf{A} assertion (\code{MODAL}, or typed
    \code{decision}/\allowbreak\code{norm}/\allowbreak\code{strategy}/\allowbreak\code{procedure_step}/\allowbreak\code{check}) that contains
    a \code{HEDGE} marker occurring \emph{before} any \code{RAT} (rationale) marker --- a hedge
    inside a \enquote{because} clause qualifies the reason, not the rule, and does not fire.
  \item[Missing element] An exception condition (\code{EXC}: \enquote{unless}, \enquote{except},
    \enquote{does not apply}, \ldots) stated by some other claim on the same topic.
  \item[Predicted knowledge type] \code{decision}; tacitness \code{relational}, or
    \code{collective} if at least half the neighborhood is typed \code{norm}.
  \item[Elicitation channel] CDM hypotheticals and boundary-case vignettes.
  \item[Question template (verbatim)] \enquote{\textquotedbl\{s1\}\textquotedbl. Describe a case
    where this did not hold, or where you decided against it. What about that case told you it was
    an exception?}
  \item[Known failure modes] (i) the \enquote{exception} found is circular, merely restating the
    judgment (e.g.\ \enquote{unless you are confident the processing is unlikely to result in a
    high risk} adds no new criterion); (ii) the hedge is statutory boilerplate with no field
    practice behind it; (iii) a trigger is stated inside the seed itself, in a form \code{EXC} does
    not match (e.g.\ \enquote{at least when}).
\end{description}

\subsection{GUARD --- named error without a detection cue}

\begin{description}
  \item[Construct and source] Automated error-monitoring and self-check habits are behavior a
    person cannot fully narrate; omission is partly a side-effect of skills becoming automatic.
    GUARD hypothesizes only the expert's own guard against an error the sources report --- it never
    asserts a \emph{learner} difficulty, because these ledgers hold expert-voiced text only, and
    experts hold just part of what learners get wrong.
  \item[Firing rule] Fires on an \textbf{A} claim typed \code{misconception}, or whose assertion
    matches \code{ERR} (\enquote{mistake}, \enquote{false positive}, \ldots) \emph{and} names a
    human agent or action (\code{GUARD_AGENT}: \enquote{technician}, \enquote{operator},
    \enquote{you}, \ldots) --- the misconception type alone no longer fires on its own.
  \item[Missing element] A detection cue (\code{DETECT}: \enquote{check}, \enquote{notice},
    \enquote{warning sign}, \ldots) stated for that error, by the seed or a neighbor.
  \item[Predicted knowledge type] \code{check} or \code{metacognition}; tacitness
    \code{automated}.
  \item[Elicitation channel] Observation or process tracing first; a retrospective probe over a
    recorded trace second, and \textbf{flagged as a weak channel} --- think-aloud is weak for
    automated, over-learned checks.
  \item[Question template (verbatim)] \enquote{\textquotedbl\{s1\}\textquotedbl. Think of the last
    time you or a colleague nearly did this or just had. What made you notice? What did you look
    at?}
  \item[Known failure modes] (i) the cue is stated in paraphrase, invisible to the \code{DETECT}
    regex; (ii) the \enquote{error} is a vendor framing (a promotional source's warning), not a
    practitioner error --- flagged live whenever the \code{PROMO} lexicon matches the span; (iii)
    the guard is organizational (a checklist or procedure), not an internalized cognitive cue.
\end{description}

\subsection{WHY --- prohibitive or contrastive prescription without a rationale}

\begin{description}
  \item[Construct and source] The rationale gap, after Szulanski's causal ambiguity: sources that
    prescribe \enquote{do this, not that} without ever saying why. Restricted to prohibitions and
    contrasts, where an omitted \enquote{why} is least likely to be simply obvious.
  \item[Firing rule] Fires on an \textbf{A} assertion with both a \code{MODAL} and a \code{CONTRA}
    marker (\enquote{not}, \enquote{never}, \enquote{instead}, \enquote{only}, \ldots), typed one
    of \code{procedure_step}/\allowbreak\code{strategy}/\allowbreak\code{check}/\allowbreak\code{decision}/\allowbreak\code{norm}, not sourced
    from a \code{standard}, and not itself matching \code{RAT}. An assertion whose only supporting
    text matches \code{AUTH} (references a regulation, statute or regulator by name) is
    \textbf{excluded as \enquote{authority}} --- logged, never emitted as a candidate at all,
    because the reason \emph{is} the rule. A descriptive \enquote{requires or does not require}
    disjunction (\code{WHY_DISJUNCTION}) also does not fire; it is not a prescription.
  \item[Missing element] A reason (\code{RAT}, or a neighbor typed \code{rationale} or the
    \code{failure_mode} the prescription guards against).
  \item[Predicted knowledge type] \code{rationale}; tacitness \code{relational}.
  \item[Elicitation channel] A targeted confirmation question.
  \item[Question template (verbatim, revised post-review)] \enquote{Sources say:
    \textquotedbl\{s1\}\textquotedbl. What would happen if someone did it differently, and in what
    situations, if any, is doing it differently acceptable?}
  \item[Known failure modes] (i) the rationale is a purpose clause (\enquote{to identify\ldots})
    that \code{RAT} does not match; (ii) the source is promotional (flagged live on a
    \code{PROMO} hit); (iii) the omitted reason is trivial safety knowledge, too obvious for any
    source to bother stating.
\end{description}

\subsection{RESULT --- PARI result interpretation}

\begin{description}
  \item[Construct and source] Added after the first adversarial review (fix F7). PARI's
    Precursor--Action--Result--Interpretation cycle: sources prescribe a test, and the expert's
    reading of the result --- the branch it selects --- is the predicted residual (the
    \enquote{A~$\to$~B~$\to$~C, but what decides B vs.\ D} gap).
  \item[Firing rule] Fires on an \textbf{A} claim typed \code{procedure_step}, \code{check},
    \code{strategy} or \code{decision}, prescriptive (\code{MODAL}, or an imperative sentence
    opening with the test verb; a safety \emph{precondition} such as \enquote{only be performed
    when\ldots} is excluded even under \code{MODAL}), whose assertion matches \code{TEST}
    (\enquote{measure}, \enquote{inspect}, \enquote{compare}, \ldots) followed within 6 tokens by
    at least one non-common content word (the object of the test).
  \item[Missing element] A result-to-interpretation mapping: what a given reading means
    (\code{INTERP}) or what value it should be compared against (\code{QUANT}).
  \item[Predicted knowledge type] \code{interpretation} plus \code{expectancy}; tacitness
    \code{relational}, plus \code{perceptual} if the test verb is \enquote{inspect} or
    \enquote{observe}.
  \item[Elicitation channel] CDM probes on a recalled case; process tracing.
  \item[Question template (verbatim)] \enquote{Sources prescribe: \textquotedbl\{s1\}\textquotedbl.
    Think of the last time you did this. What result did you get, what had you expected, what did
    that result make you do next --- and what result would have sent you down a different path?}
  \item[Known failure modes] (i) the interpretation is stated in paraphrase, missing the seed's own
    stems (RESULT shares this limitation with DISC); (ii) the test is a formality with a binary
    outcome, not a real branch point; (iii) the interpretation is instrument-given (a displayed
    pass/fail), needing no practitioner judgment at all.
\end{description}

\subsection{Retrieval-gap categories}

A retrieval gap is a separate, fourth thing, never a hidden-knowledge hypothesis
(\cref{sec:n3}). By design, each one's action is to search or verify further, not to ask a human
expert -- with one recorded exception, a code defect: \code{checks.apply_sibling} moves a record
to RG-SIBLING without clearing its \code{hypothesis} or \code{question} fields, so one GDPR v2
RG-SIBLING record (a DISC candidate) still carries both (\cref{app:defects}, row TR2).
\Cref{tab:retrieval-gaps} lists the five categories, from \code{record.categorize} and
\code{checks.apply_sibling} in \repo{gapmap/src/gapmap/record.py} and
\repo{gapmap/src/gapmap/checks.py}.

\begin{table}[htbp]
  \centering
  \small
  \begin{tabularx}{\linewidth}{@{}l X l@{}}
    \toprule
    \textbf{Category} & \textbf{What it means} & \textbf{Action} \\
    \midrule
    RG-UNK & A ledger slot that was searched for and fetched but nothing was found (an
      N2 \code{UNKNOWN} claim), or a sidecar slot marked \code{unknown}, \code{thin} or
      \code{unverified}. & Search further \\
    RG-SINGLE & A lens fired (open or partial), but only one independent source touches the
      topic ($k_{\text{topic}} \leq 1$). One voice not saying X is weak evidence that X is
      unsaid. & Search (widen sources) \\
    RG-UNVER & The lexical closure state is \code{synthetic-closed}: the missing element may be
      stated, but only in an unverified (synthetic) claim. & Verify the span \\
    RG-SIBLING & An independent sibling reconstruction of the same domain (e.g.\ GDPR v1 vs.\ v2)
      has a matching candidate that closed there; if the sibling's match is only \code{partial},
      this record downgrades from open to partial instead. & Verify (re-check against the
      sibling's sources) \\
    RG-UNDECIDED & The local judge's answer could not be parsed (malformed JSON, or it cited a
      sentence id that does not resolve). Never counted as open, never as a finding. & Re-judge
      or inspect manually \\
    \bottomrule
  \end{tabularx}
  \caption{Retrieval-gap categories and their actions. None of the five is a hidden-knowledge
    hypothesis, and by design none is routed to a human question; one exception, a code defect, is
    recorded in \cref{app:defects}, row TR2.}
  \label{tab:retrieval-gaps}
\end{table}

\paragraph{Data-wide limitations, all lenses.} All three ledgers hold only expert-voiced World A
text (documentation, procedure documents, standards); there is no work-as-done and no learner
data, so no lens can reach a learner-facing claim. Knowledge types are extractor-assigned and
directly gate DIAG and GUARD eligibility. Independence keys are coarse (3 keys in GDPR v1, 5 in
v2), which caps $k_{\text{topic}}$ and pushes most GDPR candidates to RG-SINGLE rather than HYP.
The lexicons in \repo{gapmap/src/gapmap/config.py} are English, hand-written and tuned in-sample
on these same three ledgers; they are hashed, not frozen, and must be hash-frozen at MS2 before
any gold is acquired. They have not been tested against an out-of-sample domain (\repo{docs/plans/n3-poc-lens-spec.md} \S10).
